\phantomsection
\subsection*{M3B. Vegetable Mead}
\addcontentsline{toc}{subsection}{M3B. Vegetable Mead}

\textit{Vegetable Mead deve conter pelo menos um vegetal e pode conter uma ou mais ervas ou especiarias. Os seguintes ingredientes são explicitamente incluídos como vegetais: ruibarbo, abóbora e pimentas. Transformações culinárias dos ingredientes (por exemplo, pimentas assadas) também são permitidas.}

\textbf{Impressão Geral}: Assemelha-se a um M1 Traditional Mead, porém com um caráter perceptível de vegetal e, opcionalmente, de erva ou especiaria. Uma ampla gama de resultados é possível, dependendo das escolhas do hidromel base e dos ingredientes adicionais, mas o produto final deve ser equilibrado, agradável e prazeroso. O caráter de especiaria, erva ou vegetal (do inglês {Spice, Herb and Vegetable}, SHV) não deve dominar completamente o hidromel.

\textbf{Aroma}: Os aromas do vegetal podem ser percebidos como de sutis a agressivos, refletindo o caráter dos ingredientes declarados. Os vegetais declarados devem estar perceptíveis, ainda que alguns sejam mais fortes e distintivos que outros.

Os aromas opcionais de ervas ou especiarias podem ser percebidos como de sutis a agressivos, refletindo o caráter dos ingredientes declarados. Muitos descritores são possíveis, como fragrante, aromático, pungente, herbal, picante, floral, frutado, defumado, terroso, resinoso, mentolado, ou termos que remetem a outros ingredientes (semelhante a alcaçuz, a pinho, cítrico, apimentado). As ervas e especiarias declaradas devem estar perceptíveis. Se múltiplas ervas e especiarias forem declaradas, elas não precisam estar em proporção igual, tendo em vista que algumas são mais fortes e distintivas que outras.

O caráter do mel pode variar de sutil a intenso, dependendo da força e dulçor, podendo expressar o néctar floral típico de eventuais variedades de mel declaradas. O buquê de fermentação deve ser equilibrado. Hidroméis mais fortes podem apresentar uma leve nota alcoólica picante, e os carbonatados tendem a expressar mais aroma.

Os componentes SHV e mel não precisam estar em proporção igual, mas o equilíbrio entre os ingredientes deve ser agradável e prazeroso. Os aromas provenientes de SHV devem complementar e realçar o hidromel base, não sobrepujá-lo.

\textbf{Aparência}: Limpidez de boa a brilhante; quanto maior a limpidez, mais desejável. A presença de bolhas, espuma, colarinho ou efervescência são baseados no nível de carbonatação declarado. A cor é derivada da variedade de mel utilizada. Quanto maior a vivacidade e o brilho da cor, mais desejável. A cor geralmente não é muito afetada por especiarias e ervas, embora flores e pétalas possam conferir cores de sutis a significativas.

\textbf{Sabor}: Similar ao Aroma em termos de caráter e intensidade de SHV e mel, tornando-se mais evidente nas versões mais fortes e doces. O caráter de SHV deve ser perceptível e estar em equilíbrio com o nível de dulçor do hidromel. Alguns SHVs podem trazer sabores básicos (amargo, doce, azedo, salgado, umami), além de suas características próprias.

Dulçor residual e final de acordo com o nível de dulçor declarado. A acidez deve ser equilibrada, macia ou vibrante, mas não agressiva. Taninos podem fazer um hidromel suave parecer mais seco. Características ásperas, desagradáveis ou excessivamente sulfurosas ou com notas de levedura/autólise são indesejáveis.

\textbf{Sensação na Boca}: A carbonatação é indicada pelo nível declarado, de tranquilo a espumante. O corpo geralmente aumenta com o dulçor e a força alcoólica, tipicamente de médio-leve a cheio. A percepção de álcool acompanha a força alcoólica declarada, variando de ausente a perceptível e com sensação de aquecimento. As especiarias adicionadas podem aumentar o corpo, aumentar a adstringência ou trazer notas picantes e de aquecimento; nenhum desses aspectos deve ser excessivo. Pimentas podem ser picantes, mas qualquer sensação de ardência proveniente da capsaicina deve estar em equilíbrio com o hidromel base e preservar a drinkability.

\textbf{Ingredientes}: Os mesmos ingredientes de um M1 Traditional Mead, com a adição de pelo menos um vegetal e, opcionalmente, qualquer quantidade de ervas ou especiarias.

\textbf{Comentários}: Se as especiarias forem utilizadas em conjunto com outros ingredientes, como fruta, sidra ou outros fermentáveis à base de fruta, o hidromel deve ser inscrito como M3C Fruit and Spice Mead. Se as especiarias forem utilizadas em combinação com outros tipos de ingredientes, o hidromel deve ser inscrito como M4F Experimental Mead.

\textbf{Instruções para Inscrição}: Os participantes devem especificar os níveis de dulçor, carbonatação e força alcoólica. Variedades de mel podem ser declaradas. Os participantes devem especificar o(s) vegetal(is) adicionado(s). Ervas, especiarias ou misturas de especiarias perceptíveis devem ser declaradas. Hidroméis com sabores picantes de pimenta devem especificar o nível de picância (suave, médio, forte) para auxiliar os jurados na ordenação das amostras.

\textbf{Exemplos Comerciais}: iQhilika Africa Birds Eye Chili Mead, Lyme Bay Winery Chilli Mead, Sky River Red Pepper Mead.
